% Generated draft: requires mathematical and source review.
\documentclass[11pt]{article}
\usepackage[margin=25mm]{geometry}
\usepackage{fontspec}
\setmainfont{Latin Modern Roman}
\usepackage{amsmath,amssymb,mathtools,unicode-math}
\setmathfont{Latin Modern Math}
\usepackage{enumitem,needspace,xurl,hyperref}
\setlength{\parindent}{0pt}
\setlength{\parskip}{5pt}
\setlength{\emergencystretch}{3em}
\newcommand{\N}{\mathbb{N}}
\newcommand{\Z}{\mathbb{Z}}
\newcommand{\Q}{\mathbb{Q}}
\newcommand{\R}{\mathbb{R}}
\newcommand{\C}{\mathbb{C}}
\newcommand{\F}{\mathbb{F}}
\newcommand{\A}{\mathbb{A}}
\newcommand{\PP}{\mathbb{P}}
\newcommand{\E}{\mathbb{E}}
\newcommand{\eps}{\varepsilon}
\newcommand{\dd}{\,\mathrm{d}}
\begin{document}
\section*{UnsolvedMath ID: 30005002}
\textbf{Source problem code:} catalogue entry
\subsection{Definitions and mathematical statement}
The supplied formulation is a categorical geometrization conjecture in the local setting of the cited report. Its statement depends on conventions and hypotheses omitted from the supplied passage; the following records its intended assertion while keeping these omissions explicit.

Let \(G\) be quasi-split, let \(B\subseteq G\) be a Borel subgroup, and let \(U\) be its unipotent radical. Write \(E\) for the field occurring in \(U(E)\), \(\ell\) for the coefficient prime, and \(\widehat G\) for the dual group. The supplied record does not specify the hypotheses on \(E\), the full assumptions on \(G\), or the restrictions on \(\ell\).

Let \(\psi\) be the nondegenerate character used to construct the Whittaker sheaf. The extracted passage prints \(\psi:U(E)\to\mathbb Z_\ell\), without specifying a group law or resolving whether the coefficient target has been corrupted. Consequently, the coefficient-valued character and the precise meaning of nondegeneracy remain to be fixed from the report. Subject to those conventions, its displayed construction is
\[
W_\psi=(i_1)_!\left(\operatorname{c\!\operatorname{-}Ind}_{U(E)}^{G(E)}\psi\right).
\]
Here compact induction is the construction indicated by the source, and \(i_1\) is the map used there to place the resulting object on \(\operatorname{Bun}_G\). The supplied passage does not define \(i_1\), its domain, or the sheaf-theoretic realization of compact induction.

Put
\[
\mathcal X=\operatorname{LocSys}_{\widehat G}/\mathbb Z_\ell,
\qquad
\mathcal D=D_{\mathrm{lis}}(\operatorname{Bun}_G,\mathbb Z_\ell).
\]
These denote the spectral parameter stack and the automorphic category of the report; their constructions and categorical conventions are not supplied. Let \(\operatorname{Perf}(\mathcal X)\) denote its category of perfect complexes, and write \(F*A\) for the supplied spectral action on \(\mathcal D\). The notation \(\operatorname{Coh}_{\mathrm{Nilp}}(\mathcal X)\) means the source's category of coherent complexes satisfying the Arinkin--Gaitsgory singular-support condition
\[
\operatorname{SS}(F)\subseteq\operatorname{Nilp},
\]
where \(\operatorname{Nilp}\) is its nilpotent cone in the space governing singular support. The source identifies the zero-section support condition with the perfect subcategory. The precise derived model, support space, and conventions for the coherent category are not specified in the record.

\textbf{Conjectural assertion.} With all these objects interpreted in the conventions of the report, the functor
\[
\Phi_0:\operatorname{Perf}(\mathcal X)\longrightarrow\mathcal D,
\qquad F\longmapsto F*W_\psi,
\]
extends along the inclusion of perfect complexes to an equivalence
\[
\Phi:\operatorname{Coh}_{\mathrm{Nilp}}(\mathcal X)
\xrightarrow{\ \sim\ }\mathcal D
\]
compatible with the spectral action. The source further states that the nilpotent singular-support condition becomes automatic after inverting \(\ell\), whereas it is needed integrally. This qualification is retained as a source assertion; its precise scope depends on the omitted categorical conventions. These omissions prevent treating the displayed formulation as a fully specified standalone conjecture.
\subsection{Short English statement}
The conjecture says that acting on a Whittaker sheaf should identify the integral spectral category with nilpotent singular support with the automorphic category on \(\operatorname{Bun}_G\), compatibly with the spectral action. The source says the support restriction becomes automatic after inverting \(\ell\). Field hypotheses, character coefficients, and category conventions require clarification. The supplied dated curation labels the problem partially solved while describing the full equivalence as conjectural; this draft does not independently verify either assessment.
\subsection{Sources}
\begin{itemize}
\item Local UnsolvedMath registry, permanent ID \(30005002\), exact source problem code \texttt{catalogue entry}, title \emph{Integral Geometrization with Nilpotent Singular Support}; supplied complete source record. No individual proposer is supplied; the proposed year is recorded as 2022.
\item Supplied report citation: \emph{Non-Archimedean Geometry and Applications (2023)}, Oberwolfach Reports, pp.\ 231--302, DOI 10.4171/owr/2022/5. \url{https://doi.org/10.4171/owr/2022/5}. The literature-source list labels this same reference \emph{Oberwolfach Report 5/2022}; the extracted passage is headed ``6. The geometrization conjecture.''
\item Fargues--Scholze, \emph{Geometrization of the local Langlands correspondence}. Supplied literature reference: \url{https://arxiv.org/abs/2102.13459}.
\item Literature assessment dated exactly \texttt{2026-08-22}: supplied dated curation, not independently verified. It labels the problem partially solved, credits the Fargues--Scholze program with major foundational and spectral-action results, and describes the stated categorical equivalence with nilpotent singular support as still conjectural in this generality. The background explicitly requests independent verification of time-sensitive status and cited claims.
\item The registry labels the statement \texttt{corrected\_verified} and reports reconstruction and checking against the official Oberwolfach report. This is supplied provenance, not an independent verification performed for this draft.
\end{itemize}
\end{document}
